\documentclass[12pt, a4paper]{article}
\usepackage[utf8]{inputenc}
\usepackage[margin=1in]{geometry}
\usepackage{xcolor}
\usepackage{titlesec}
\usepackage{enumitem}
\usepackage{hyperref}

% Basic Color Formatting
\titleformat{\section}{\large\bfseries\color{blue}}{}{0em}{}[\titlerule]
\titleformat{\subsection}{\normalsize\bfseries}{}{0em}{}

\title{The Four Pillars of Object-Oriented Programming: A Crash Course}
\author{Computer Science Education Series}
\date{\today}

\begin{document}

\maketitle

\section{Introduction}
Object-Oriented Programming (OOP) is a paradigm based on the concept of "objects," which can contain data and code[cite: 4, 5].

\section{1. Abstraction}
Abstraction hides implementation details while showing functionality[cite: 4, 5].

\section{2. Encapsulation}
Encapsulation binds together code and data, keeping both safe from outside interference[cite: 4, 5].

\section{3. Inheritance}
Inheritance allows one object to acquire the properties and behaviors of a parent object[cite: 4, 5].

\section{4. Polymorphism}
Polymorphism allows objects of different types to be treated as objects of a common base type[cite: 4, 5].

\begin{thebibliography}{9}
\bibitem{grady2007} 
Booch, G., et al. (2007). \textit{Object-Oriented Analysis and Design with Applications}.
\end{thebibliography}

\end{document}
